最大子数组问题（Maximum Subarray）：给定一个整数数组，
找出一段连续子数组，使得它的元素之和最大。这个问题源于股票买卖
—— 把每天的股价差当作数组元素，找到「哪天买、哪天卖」最划算。

你在项目中实现了\textbf{两种}解法，很好地展示了从朴素到高效的演进。

\subsection{解法一：暴力枚举所有可能的子数组}

对每一对下标 \((i, j)\) 都计算 \(\sum A[i..j]\)，取最大。
共 \(O(n^2)\) 个子数组，每个求和 \(O(1)\)（因为用 \texttt{A[j] - A[i]}
的差分思路，实际只计算了两端差值）。

\lstinputlisting[style=cstyle,
    caption=algo/src/findmaximumsubarray/1.c — 暴力枚举法]
{../algo/src/findmaximumsubarray/1.c}

这段代码的思路：

\begin{itemize}
    \item 对所有 \(i < j\)，计算 \texttt{A[j] - A[i]}，代表「从 i 买入、j 卖出」的收益
    \item 顺便记录是哪一对 \((i, j)\) 产生了最大收益
    \item 时间复杂度 \(O(n^2)\)，空间 \(O(n)\) 来存所有对
\end{itemize}

对小规模数据（比如 \(n = 17\)）这样写完全没问题，结果正确，
但当 \(n\) 达到几万或更大时，二次方的开销就难以接受了。

\subsection{解法二：分治法}

把数组从中间 \texttt{mid} 切开。最大子数组必然落在以下三种情况之一：

\begin{enumerate}
    \item 完全在左半部分 \texttt{A[low..mid]}
    \item 完全在右半部分 \texttt{A[mid+1..high]}
    \item 跨越 \texttt{mid}：一部分在左、一部分在右
\end{enumerate}

三种情况都递归处理，最后取三者中最大的。

\subsubsection{跨中点的最大子数组 — \texttt{findmaxcrossingsubarray}}

\lstinputlisting[style=cstyle, firstline=1, lastline=35,
    caption=algo/src/findmaximumsubarray/2.c — 跨中点的查找]
{../algo/src/findmaximumsubarray/2.c}

逻辑：

\begin{itemize}
    \item 从 \texttt{mid} 向左遍历，累加求和并更新最大值和对应位置 \texttt{max\_left}
    \item 从 \texttt{mid+1} 向右遍历，同样记录 \texttt{max\_right}
    \item 结果是 \texttt{(max\_left, max\_right, leftsum + rightsum)}
    \item 这一步是 \(O(n)\) 的，因为只从中间向两边各扫描一次
\end{itemize}

注意你设 \texttt{leftsum} 和 \texttt{rightsum} 的初值为 \texttt{-1000000}。
这个\textbf{哨兵值}很重要——因为数组元素可以是负数，必须让首次比较
「哪怕是一个负数的子数组」也能被选中。如果设为 0，当所有数都是负数
时算法就会返回错误的空数组。

\subsubsection{主递归函数}

\lstinputlisting[style=cstyle, firstline=36, lastline=61,
    caption=algo/src/findmaximumsubarray/2.c — findmaximumsubarray 递归函数]
{../algo/src/findmaximumsubarray/2.c}

\subsubsection{输入预处理和输出结果}

\lstinputlisting[style=cstyle, firstline=62, lastline=75,
    caption=algo/src/findmaximumsubarray/2.c — main，把股价转成差分数组]
{../algo/src/findmaximumsubarray/2.c}

\subsection{时间复杂度对比}

\begin{center}
\begin{tabular}{llll}
解法 & 时间 & 空间 & 备注 \\
暴力（解法一） & \(O(n^2)\) & \(O(n)\) & 代码最短，适合调试 \\
分治（解法二） & \(O(n \log n)\) & \(O(\log n)\) 递归栈 & 大规模数据 \\
\end{tabular}
\end{center}

你在项目中保留了两种实现，正好可以对照运行来验证正确性——
这是学习算法的一种很实用的方式：\textbf{写一个简单但正确的版本，
再写一个快的版本，让两者对拍}。
